% =========================================================================
% Chapter 5: Conclusion and Future Work
%
% Student Guide:
% - Sums up what you found, what it adds, its limits and what could come next.
% - Do not add new data or new arguments here.
% =========================================================================

\chapter{Conclusion}
\label{ch:conclusion}

The Conclusion chapter is the final section of a thesis and serves as the place where the researcher ties together all elements of the study. It should not introduce new data or arguments but instead provide a reflective closure that synthesizes the research problem, findings, contributions, limitations, and future directions. Students should approach this chapter as both a summary and a forward-looking statement, ensuring that it is concise, cohesive, and academically rigorous. The goal is to leave the reader with a clear understanding of the research impact and its relevance to the broader field of Computer Science and Engineering.

\section{Restating the Research Problem and Objectives}
\label{sec:restating-problem-objectives}

Begin by revisiting the central research problem and the objectives that guided the study. Present them as a reflective summary that frames the research in light of what has been accomplished. This establishes continuity and prepares the reader to understand the significance of the findings.

\section{Summarizing Key Findings}
\label{sec:summarizing-findings}

Provide a concise synthesis of the major findings. Highlight how the research questions were answered or how the objectives were met. Focus on the broader implications of the results rather than technical detail or raw data.

\section{Highlighting Contributions}
\label{sec:ch5-highlighting-contributions}

Clearly articulate the contributions of the thesis. These may be theoretical, methodological, or practical. State explicitly how the work adds value to the academic community and industry, and distinguish it from existing literature.

\section{Discussing Limitations}
\label{sec:discussing-limitations}

Include a paragraph that acknowledges the limitations of the study. Discuss constraints such as scope, resources, or assumptions. Maintain a balanced tone that recognizes limitations without undermining the overall contribution.

\section{Suggesting Future Work}
\label{sec:suggesting-future-work}

Offer realistic and specific directions for future research. In Computer Science and Engineering, where technology evolves rapidly, suggestions might include scaling the proposed solution to larger datasets, testing in real-world environments, or integrating with emerging technologies. The future work section should inspire subsequent researchers while remaining grounded in the scope of the thesis.
